\documentclass[11pt,reqno]{amsart}
\usepackage[T1]{fontenc}
\usepackage[utf8]{inputenc}
\usepackage{lmodern}
\usepackage{amsmath,amssymb,amsfonts,amsthm,mathtools,mathrsfs}
\usepackage{enumitem}
\usepackage{microtype}
\usepackage[colorlinks=true,linkcolor=blue,citecolor=blue,urlcolor=blue]{hyperref}
\usepackage[nameinlink,capitalize]{cleveref}
\allowdisplaybreaks
\setlength{\emergencystretch}{2em}

\newtheorem{theorem}{Theorem}[section]
\newtheorem{proposition}[theorem]{Proposition}
\newtheorem{lemma}[theorem]{Lemma}
\newtheorem{corollary}[theorem]{Corollary}
\theoremstyle{definition}
\newtheorem{definition}[theorem]{Definition}
\newtheorem{assumption}[theorem]{Assumption}
\newtheorem{example}[theorem]{Example}
\theoremstyle{remark}
\newtheorem{remark}[theorem]{Remark}

\newcommand{\A}{\mathcal A}
\newcommand{\B}{\mathcal B}
\newcommand{\C}{\mathbb C}
\newcommand{\E}{\mathbb E}
\newcommand{\R}{\mathbb R}
\newcommand{\T}{\mathbb T}
\newcommand{\Piop}{\Pi}
\newcommand{\Id}{\mathrm{Id}}
\newcommand{\tr}{\operatorname{tr}}
\newcommand{\spec}{\operatorname{spec}}
\newcommand{\dist}{\operatorname{dist}}
\newcommand{\norm}[1]{\left\lVert #1\right\rVert}
\newcommand{\abs}[1]{\left\lvert #1\right\rvert}
\newcommand{\pair}[2]{\left\langle #1,#2\right\rangle}

\title[Pressure, diffusion, and full-frequency response]{Pressure, Physical Diffusion, and Uniform Full-Frequency Suspension Response for Moving Collision Systems}
\author{Qian Qi}
\address{Peking University, Beijing 100871, China}
\email{qiqian@pku.edu.cn}
\subjclass[2020]{Primary 37A50, 37C30; Secondary 47A55, 60F17}
\keywords{pressure, suspension flow, Green--Kubo formula, diffusion, high-frequency resolvent, moving billiards}
\date{August 29, 2026}

\begin{document}

\begin{abstract}
We convert all-order collision response into physical-time pressure, diffusion,
and full-frequency suspension resolvents.  A finite Banach bundle is stabilized
inside one direct-sum operator space.  Joint displacement--roof twists then
produce differentiable pressure and an implicit zero-pressure root.  Its
Hessian gives physical covariance, while an oriented renewal identity uses
distinct entry and exit roof-cell operators.

The diffusion tensor is defined by the symmetrized Green--Kubo series and is
proved equal to the pressure Hessian, the asymptotic variance, and the Gordin
martingale bracket.  Symmetry and positive semidefiniteness are therefore
conclusions.  Uniform positivity yields a quantitative square-root modulus and
a uniformly elliptic macroscopic coefficient field.

For the nonconjugate four-branch moving collision family of the companion
paper, branchwise roofs and displacements make the pressure explicit.  A
Diophantine roof vector gives a polynomial lower bound on the scalar phase
defect at every nonzero frequency; the quotient by constants contracts
uniformly.  This yields a full-frequency resolvent theorem and parameter
derivatives of every prescribed finite order on the same moving family.  For
analytic no-eclipse open dispersing billiards, the fixed symbolic coding and a
verifiable temporal-shear condition give a uniform Dolgopyat estimate and
moving-family resolvent derivatives.  Finally, we prove the exact finite-
horizon interval for the triangular periodic Lorentz gas:
$\sqrt3/4<r<1/2$, together with a quantitative deformation margin.
\end{abstract}

\maketitle
\tableofcontents

\section{Introduction}

A collision map is naturally parameterized by collision count, while a slow
limit evolves in physical time.  Passing between the two requires a joint
spatial and roof twist, an implicit physical-time pressure root, and a renewal
formula with correctly oriented incomplete roof segments.  At high frequency,
a family theorem needs a common domain and a quantitative mechanism preventing
the twisted leading eigenvalue from returning too close to one.

The companion response paper supplies two actual moving systems.  The first is
an explicit symplectic Markov collision map with all-order response and exact
innovations.  The second is an analytic no-eclipse open dispersing billiard
pulled to one fixed symbolic space.  We prove full-frequency theorems for both.
For the Markov system every estimate is explicit.  For the open billiard, a
geometric temporal shear is fed into the Dolgopyat cancellation mechanism on
the fixed symbolic space.

The perturbative ingredients go back to Kato and Keller--Liverani
\cite{Kato1995,KellerLiverani1999}.  Pressure and symbolic suspension theory are
classical \cite{ParryPollicott1990}.  High-frequency estimates originate with
Dolgopyat \cite{Dolgopyat1998}; the open-billiard version uses the transfer
operator technology of Stoyanov \cite{Stoyanov2011}.  Our contribution is the
moving-family assembly, explicit phase separation, parameter derivatives of
the full resolvent, and the exact identification of the physical covariance.

\section{A fixed ambient operator space}

Let $a$ range over a compact parameter manifold $\A$.  Let $\B_a$ be a Banach
bundle with finitely many charts $U_\alpha$, model spaces $\B_\alpha$, and
isomorphisms
\[
 G_{\alpha,a}:\B_a\longrightarrow\B_\alpha
\]
whose transition maps and inverses are uniformly bounded and $C^N$ in $a$.
Choose smooth functions $\psi_\alpha$ subordinate to the atlas such that
\begin{equation}\label{eq:quadratic-partition}
 \sum_{\alpha=1}^M\psi_\alpha(a)^2=1.
\end{equation}
Put $\B_*=igoplus_\alpha\B_\alpha$ and define
\begin{align}
 J_ah&=(\psi_\alpha(a)G_{\alpha,a}h)_\alpha,\label{eq:J}\\
 R_a(v_\alpha)_\alpha&=
 \sum_\alpha\psi_\alpha(a)G_{\alpha,a}^{-1}v_\alpha.\label{eq:R}
\end{align}

\begin{theorem}[Finite-atlas stabilization]\label{thm:stabilization}
One has $R_aJ_a=\Id_{\B_a}$.  Hence $P_a=J_aR_a$ is a uniformly bounded
projection onto $E_a=J_a\B_a\subset\B_*$.  If $L_a:\B_a\to\B_a$, then
\[
 \widehat L_a=J_aL_aR_a\in\mathcal L(\B_*)
\]
has the same parameter regularity.  Riesz projections around an isolated
nonzero spectral value agree with the fibre projections after restriction to
$E_a$.
\end{theorem}

\begin{proof}
Equation \eqref{eq:quadratic-partition} gives
$R_aJ_ah=\sum_\alpha\psi_\alpha^2h=h$.  Thus $P_a^2=P_a$ and its range is
$E_a$.  Finiteness of the atlas gives all uniform bounds and derivatives.  The
stabilized operator vanishes on the chosen complement, so the additional
spectral point zero is separated from a leading eigenvalue near one.
\end{proof}

\section{Joint pressure and the physical-time root}

Let $\kappa_a:M\to\R^d$ be a displacement and
$\tau_a:M\to(0,\infty)$ a roof.  Define
\begin{equation}\label{eq:twist}
 L_{a,q,s}h=L_a\left(e^{q\cdot\kappa_a-s\tau_a}h\right).
\end{equation}
Assume the all-order graded letters supplied by the companion paper and a
uniform contour around one simple eigenvalue $\lambda(a,q,s)$ near one.  Put
\[
 \mathscr P(a,q,s)=\log\lambda(a,q,s),
 \qquad \mathscr P(a,0,0)=0.
\]

\begin{theorem}[Pressure jet]\label{thm:pressure}
For every prescribed finite order $N$, the eigenvalue, Riesz projection,
reduced resolvent, and pressure are $C^N$ in the real variables $(a,q,s)$.
Every derivative is a finite sum of well-typed reduced-resolvent words whose
total parameter order is at most $N$.
\end{theorem}

\begin{proof}
On the fixed space,
\[
 \Piop_{a,q,s}={1\over2\pi i}\int_\Gamma
 (z-\widehat L_{a,q,s})^{-1}\,dz.
\]
Repeated differentiation uses
$D(z-L)^{-1}=(z-L)^{-1}(DL)(z-L)^{-1}$.  The graded response theorem makes
every word bounded, and uniform contour separation permits differentiation
under the integral.  A normalized left--right eigenpair gives the eigenvalue
and its logarithm.
\end{proof}

Let $\Lambda_a(q)$ solve
\begin{equation}\label{eq:root}
 \mathscr P(a,q,\Lambda_a(q))=0.
\end{equation}
Since
\[
 \mathscr P_s(a,0,0)=-\bar\tau_a,
 \qquad \bar\tau_a=\int\tau_a\,d\mu_a>0,
\]
the implicit-function theorem applies.

\begin{theorem}[Physical drift and covariance]\label{thm:physical}
The physical root is $C^N$.  If the displacement is centered, then
$D_q\Lambda_a(0)=0$ and
\begin{equation}\label{eq:Sigma}
 \Sigma_{ij}(a)=\partial_{q_iq_j}\Lambda_a(0)
 =\frac{\mathscr P_{q_iq_j}(a,0,0)}{\bar\tau_a}.
\end{equation}
In particular,
\begin{equation}\label{eq:Sigma-response}
 \partial_a\Sigma_{ij}
 =\frac{\mathscr P_{aq_iq_j}}{\bar\tau}
 -\frac{\mathscr P_{q_iq_j}\partial_a\bar\tau}{\bar\tau^2}.
\end{equation}
The physical diffusion matrix is $D^{\rm phys}=\Sigma/2$.
\end{theorem}

\begin{proof}
Differentiate \eqref{eq:root}.  Centering removes the first-gradient cross
terms.  Twice differentiating in $q$ gives
$0=\mathscr P_{q_iq_j}+\mathscr P_s\Lambda_{q_iq_j}$, which yields
\eqref{eq:Sigma}.  Differentiation in $a$ gives
\eqref{eq:Sigma-response}.
\end{proof}

\section{Oriented suspension renewal}

For a flow observable $G$, define the entry operator
\begin{equation}\label{eq:entry}
 \mathcal B_{G,a,z}(x)=
 \int_0^{\tau_a(x)}e^{-z(\tau_a(x)-u)}G(x,u)\,du.
\end{equation}
For $F$, define the exit functional
\begin{equation}\label{eq:exit}
 \ell_{F,a,z}(h)=
 \int_M\int_0^{\tau_a(x)}e^{-zu}F(x,u)h(x)\,du\,d\mu_a(x).
\end{equation}
Let $L_{a,z}=L_a(e^{-z\tau_a}\cdot)$.

\begin{theorem}[Oriented renewal identity]\label{thm:renewal}
The Laplace transform of a centered suspension correlation is
\begin{equation}\label{eq:renewal}
 \widehat C_{F,G}(a,z)=H_{a,z}^{\rm local}(F,G)
 +\ell_{F,a,z}(I-L_{a,z})^{-1}\mathcal B_{G,a,z},
\end{equation}
where the local term contains trajectories with no complete return between the
observations.  Near $z=0$, the centered expression has the same parameter
regularity as the pressure jet.
\end{theorem}

\begin{proof}
Decompose a trajectory by the number of complete returns.  The first
incomplete segment is encoded by \eqref{eq:entry}, the last by
\eqref{eq:exit}, and the complete returns by powers of $L_{a,z}$.  Summing the
geometric series gives \eqref{eq:renewal}.  Near zero split
\[
 (I-L_{a,z})^{-1}=
 {\Piop_{a,z}\over1-\lambda(a,0,z)}+R_{a,z}^{\perp}.
\]
Centering cancels the pole and
$\partial_z\lambda(a,0,0)=-\bar\tau_a$ is uniformly nonzero.
\end{proof}

\section{Symmetric Green--Kubo covariance}

For centered $\kappa_a$ define
\begin{equation}\label{eq:GK}
 C_{ij}^{\rm coll}(a)=
 \int\kappa_{a,i}\kappa_{a,j}\,d\mu_a
 +\sum_{n\ge1}\int\left(
 \kappa_{a,i}\kappa_{a,j}\circ T_a^n+
 \kappa_{a,j}\kappa_{a,i}\circ T_a^n
 \right)d\mu_a.
\end{equation}
Assume absolute summability, which is automatic in the actual models below.

\begin{theorem}[Pressure--variance--bracket identity]\label{thm:GK}
One has
\begin{equation}\label{eq:GK-identities}
 \mathscr P_{q_iq_j}(a,0,0)=C_{ij}^{\rm coll}(a)
 =\int m_{a,i}m_{a,j}\,d\mu_a,
\end{equation}
where $m_a$ is the Gordin martingale increment.  Equivalently,
\[
 \xi^TC^{\rm coll}(a)\xi
 =\lim_{n\to\infty}{1\over n}
 \int\left(\sum_{k=0}^{n-1}
 \xi\cdot\kappa_a\circ T_a^k\right)^2d\mu_a\ge0.
\]
Thus $C^{\rm coll}$ and $\Sigma=C^{\rm coll}/\bar\tau$ are symmetric positive
semidefinite.  Their kernel is precisely the coboundary subspace in the
declared spectral class.
\end{theorem}

\begin{proof}
Twice differentiate the leading eigenvalue equation in $q$.  The direct term
is the instantaneous covariance and the two reduced-resolvent insertions
expand into the two ordered correlation series in \eqref{eq:GK}.  Expanding
the square of partial sums and using absolute summability gives the variance
limit.  Solving the Poisson equation and expanding the martingale square gives
the bracket formula.  A zero bracket is equivalent to a zero martingale part,
hence to a coboundary.
\end{proof}

\begin{corollary}[Quantitative square-root modulus]\label{cor:sqrt}
If $\Sigma(a)\ge\lambda I$ on a compact parameter set, then
\[
 \norm{\Sigma(a)^{1/2}-\Sigma(a')^{1/2}}
 \le {1\over2\sqrt\lambda}
 \norm{\Sigma(a)-\Sigma(a')}.
\]
\end{corollary}

\begin{proof}
The difference of the square roots solves a Sylvester equation whose two
spectra are bounded below by $\sqrt\lambda$.
\end{proof}

The antisymmetric iterated-correlation tensor is recorded separately as an
area anomaly; it is not inserted into the diffusion matrix.

\section{The explicit moving-collision pressure model}

Use the four-branch family of the companion paper.  Choose branch
displacements $c_i\in\R^d$ and positive roofs $\tau_i(a)$.  The joint pressure
is
\begin{equation}\label{eq:explicit-pressure}
 \mathscr P(a,q,s)=
 \log\left(\sum_{i=1}^4w_i(a)e^{q\cdot c_i-s\tau_i(a)}\right).
\end{equation}
Let
\[
 \bar c(a)=\sum_iw_i(a)c_i,
 \qquad \widetilde c_i(a)=c_i-\bar c(a).
\]

\begin{theorem}[Actual physical coefficient package]\label{thm:actual}
For the moving collision system:
\begin{enumerate}[label=(\roman*)]
\item the pressure and physical root are analytic;
\item all positive-time correlations of branchwise centered displacement
vanish and
\[
 C^{\rm coll}(a)=\sum_iw_i(a)
 \widetilde c_i(a)\otimes\widetilde c_i(a);
\]
\item $\Sigma(a)=C^{\rm coll}(a)/\bar\tau(a)$;
\item if the centered branch vectors span $\R^d$ uniformly, then
$\Sigma(a)\ge\lambda I$;
\item every coefficient derivative appearing in the downstream
homogenization is an explicit derivative of \eqref{eq:explicit-pressure}.
\end{enumerate}
\end{theorem}

\begin{proof}
The constant function is an eigenfunction of the twisted full-branch transfer
operator with eigenvalue equal to the sum in \eqref{eq:explicit-pressure}.
The quotient by constants is uniformly contracting, so this eigenvalue is
simple near the origin.  Exact innovations make successive branch labels
conditionally independent with probabilities $w_i(a)$; at frozen $a$ they are
i.i.d.  The covariance is therefore the one-step variance.  Uniform spanning
and compactness give the lower eigenvalue bound.
\end{proof}

\section{A full-frequency theorem for the moving collision family}

For $z=\sigma+ib$ define
\begin{equation}\label{eq:twisted-full-branch}
 (\mathcal L_{a,z}h)(y)=
 \sum_{i=1}^4w_i(a)e^{-z\tau_i(a)}
 h(s_{i-1}(a)+w_i(a)y).
\end{equation}
The constant subspace is invariant, with scalar eigenvalue
\begin{equation}\label{eq:lambda-z}
 \lambda_a(z)=\sum_iw_i(a)e^{-z\tau_i(a)}.
\end{equation}
On $W^{r,1}/\C1$, use the derivative quotient norm.

\begin{lemma}[Uniform quotient contraction]\label{lem:quotient}
For $\operatorname{Re}z\ge0$, the operator induced by
$\mathcal L_{a,z}$ on $W^{r,1}/\C1$ has norm at most $\rho<1$, uniformly in
$a$ and $z$.
\end{lemma}

\begin{proof}
The phase factors have modulus at most one.  Differentiating in $y$ and
changing variables on each branch gives
\[
 \norm{(\mathcal L_{a,z}h)'}_1
 \le\rho\norm{h'}_1.
\]
The same estimate at higher derivative levels gives the quotient bound.
\end{proof}

\begin{assumption}[Diophantine roof vector]\label{ass:diophantine}
There are $c_0>0$ and $\nu\ge1$ such that for all $\abs b\ge1$,
\begin{equation}\label{eq:diophantine}
 \inf_{k\in\mathbb Z^4}
 \max_{1\le i\le4}\abs{b\tau_i-2\pi k_i}
 \ge c_0(1+\abs b)^{-\nu}.
\end{equation}
\end{assumption}
Such roof vectors form a full-measure class for every exponent above the
simultaneous approximation threshold.

\begin{lemma}[Scalar phase separation]\label{lem:phase}
Under \cref{ass:diophantine}, with branch weights bounded below by $w_*>0$,
\begin{equation}\label{eq:phase-gap}
 \abs{1-\lambda_a(ib)}
 \ge c_1(1+\abs b)^{-2\nu},
 \qquad \abs b\ge1,
\end{equation}
uniformly in $a$.
\end{lemma}

\begin{proof}
For phases $\theta_i=b\tau_i$,
\[
 1-\abs{\textstyle\sum_iw_ie^{-i\theta_i}}^2
 =2\sum_{i<j}w_iw_j(1-\cos(\theta_i-\theta_j)).
\]
If the weighted sum is not close to one, the conclusion is immediate.  If it
is close to one, all phases must lie in one short arc and their common phase
must also be close to zero modulo $2\pi$.  Uniform convexity of the unit circle
then gives
\[
 \abs{1-\textstyle\sum_iw_ie^{-i\theta_i}}
 \ge c\inf_{k\in\mathbb Z^4}
       \max_i\abs{\theta_i-2\pi k_i}^2.
\]
Apply \eqref{eq:diophantine}.
\end{proof}

\begin{theorem}[Full-frequency moving-family resolvent]\label{thm:full-frequency}
Assume \cref{ass:diophantine}.  For
$0\le\operatorname{Re}z\le\sigma_0$ and $\abs{\operatorname{Im}z}\ge1$,
\begin{equation}\label{eq:resolvent-bound}
 \norm{(I-\mathcal L_{a,z})^{-1}}_{W^{r,1}\to W^{r,1}}
 \le C_r(1+\abs{\operatorname{Im}z})^{2\nu}.
\end{equation}
For every prescribed parameter order $N$,
\begin{equation}\label{eq:resolvent-derivative}
 \norm{\partial_a^N(I-\mathcal L_{a,z})^{-1}}
 _{W^{r+N,1}\to W^{r,1}}
 \le C_{r,N}(1+\abs{\operatorname{Im}z})^{2\nu(N+1)+N}.
\end{equation}
Together with the low-frequency spectral decomposition, these estimates give
the full-frequency suspension response through order $N$.
\end{theorem}

\begin{proof}
Relative to constants and the quotient, the operator is upper triangular.  The
constant block is $\lambda_a(z)$ and the quotient block has norm at most
$\rho$.  The inverse block formula, \cref{lem:quotient,lem:phase}, gives
\eqref{eq:resolvent-bound}; the off-diagonal block is uniformly bounded by the
same adapted Sobolev norm.  Differentiating the inverse produces sums of words
\[
 R(\partial_a^{k_1}\mathcal L)R\cdots
 (\partial_a^{k_j}\mathcal L)R,
 \qquad k_1+\cdots+k_j=N.
\]
A parameter derivative of the roof multiplier costs at most one factor of
$1+\abs z$ per derivative, and a derivative of the inverse branch loses the
same number of Sobolev levels as in the companion paper.  Combining these
bounds with $j+1\le N+1$ resolvents gives
\eqref{eq:resolvent-derivative}.
\end{proof}

\section{Uniform high-frequency response for moving open billiards}

Use the analytic no-eclipse open billiard and fixed symbolic coding from the
companion paper.  Let $h_1,h_2$ be two symbolic inverse branches of a common
return map and let $\tau_N$ be the $N$-return roof sum.

\begin{assumption}[Temporal shear]\label{ass:shear}
There are a cylinder $U$, two inverse branches, a unit symbolic tangent
direction $v$, and $\kappa>0$ such that
\begin{equation}\label{eq:shear}
 \abs{D_v(\tau_N\circ h_1-\tau_N\circ h_2)(\xi)}\ge\kappa,
 \qquad \xi\in U,\quad a\in\A.
\end{equation}
\end{assumption}

\begin{proposition}[Concrete shear for three asymmetric disks]\label{prop:disk-shear}
For three circular obstacles whose centers are noncollinear and whose radii and
separations satisfy no eclipse, an open dense set of configurations satisfies
\cref{ass:shear}.  The condition persists under small analytic motions of one
obstacle.
\end{proposition}

\begin{proof}
Take the two return words from obstacle one through obstacle two and through
obstacle three.  Differentiating their length sums with respect to the starting
boundary coordinate on obstacle one gives the difference of the tangential
projections of the two outgoing and incoming unit directions.  At a normal
impact this difference is nonzero unless the two relevant center directions
are symmetric about the normal.  Noncollinearity and an arbitrarily small
asymmetric displacement avoid that codimension-one relation.  Compactness of a
smaller cylinder gives a positive lower bound, and all quantities vary
continuously with the obstacle family.
\end{proof}

\begin{theorem}[Open-billiard moving-family Dolgopyat response]\label{thm:open-HF}
Under the analytic no-eclipse hypotheses and \cref{ass:shear}, there are
$b_0,C,\eta>0$ such that the normalized symbolic suspension operator satisfies
\[
 \norm{(I-\mathscr L_{a,ib})^{-1}}_{C^\alpha\to C^\alpha}
 \le C\abs b^\eta,
 \qquad \abs b\ge b_0,\quad a\in\A.
\]
For every finite $N$, all parameter derivatives of the resolvent exist and
obey a polynomial bound whose exponent depends only on $N$ and the uniform
Dolgopyat constants.  Hence analytic physical observables have a uniform
full-frequency moving-scatterer suspension response.
\end{theorem}

\begin{proof}
The fixed symbolic coding turns the moving billiard into a $C^\omega$ family
of Hölder roof and Jacobian potentials.  The temporal shear
\eqref{eq:shear} gives phase separation between the two inverse branches on
$U$.  The Dolgopyat cone argument then yields a contraction of a logarithmic
iterate in the $b$-weighted Hölder norm, uniformly in $a$; see
\cite{Dolgopyat1998,Stoyanov2011}.  Summing the contracted blocks gives the
polynomial resolvent bound.  The shear, distortion, and cone margins are open
and uniform on the compact family.  Parameter differentiation is performed on
the fixed Hölder space; the resolvent identity gives finite words of resolvents
and analytic potential derivatives, hence the asserted polynomial bounds.
\end{proof}

\section{The triangular Lorentz finite-horizon interval}

Let
\[
 \Lambda=\mathbb Z(1,0)+\mathbb Z(1/2,\sqrt3/2)
\]
and place a closed disk of radius $r$ at each lattice point.

\begin{theorem}[Exact corridor threshold]\label{thm:triangle}
The disks are disjoint and the periodic Lorentz gas has finite horizon for
\begin{equation}\label{eq:triangle-window}
 {\sqrt3\over4}<r<{1\over2}.
\end{equation}
If each disk is replaced by a strictly convex $C^2$ obstacle containing the
concentric disk of radius $r-\delta$, with
$r-\delta>\sqrt3/4$, and remaining inside the disk of radius $r+\delta<1/2$,
then finite horizon, nonoverlap, and a positive curvature margin persist.
\end{theorem}

\begin{proof}
For a primitive lattice vector $v$, adjacent lattice lines parallel to $v$
are separated by
\[
 {\operatorname{area}(\R^2/\Lambda)\over\abs v}
 ={\sqrt3/2\over\abs v}.
\]
The largest separation occurs for the shortest primitive vectors, for which it
is $\sqrt3/2$.  A straight corridor can exist only if one such strip has width
larger than $2r$.  Thus $r>\sqrt3/4$ blocks every primitive direction.  Since
the unit tangent bundle modulo the lattice is compact, absence of a straight
corridor implies a finite maximum free-flight length.  The upper bound
$r<1/2$ is exactly nonoverlap for nearest neighbours.  The same line-strip
argument applied to the contained disks proves the deformation statement;
strict curvature and nonoverlap are preserved by the stated margins.
\end{proof}

\section{Macroscopic coefficient fields}

Let $\Theta:\R^m\times\R^m\to\A$ and $B(x,p)$ be globally
$C^{N,\alpha}$ with bounded derivatives.  Define
\[
 A(x,p)={1\over2}B(x,p)\Sigma(\Theta(x,p))B(x,p)^T.
\]

\begin{theorem}[Coefficient lift and ellipticity]\label{thm:lift}
All pressure, drift, roof, covariance, area, and response fields pulled back by
$\Theta$ are $C^{N,\alpha}$.  If
\[
 \Sigma(a)\ge\lambda_\Sigma I,
 \qquad B(x,p)B(x,p)^T\ge\lambda_B^2I,
\]
then
\[
 A(x,p)\ge{1\over2}\lambda_\Sigma\lambda_B^2I.
\]
\end{theorem}

\begin{proof}
The fixed operator realization gives parameter regularity before composition.
The finite-dimensional chain rule gives the pulled-back fields.  For $z\in
\R^m$,
\[
 z^TAz={1\over2}(B^Tz)^T\Sigma(B^Tz)
 \ge{1\over2}\lambda_\Sigma\lambda_B^2\abs z^2.
\]
\end{proof}

\section{Conclusion}

The collision-to-physical-time interface is now positive at both low and high
frequency.  The symmetrized Green--Kubo tensor is identified with pressure,
variance, and martingale bracket.  The moving Markov collision family has an
explicit Diophantine full-frequency theorem with parameter derivatives, while
the open dispersing billiard has a uniform moving-family Dolgopyat theorem on
its fixed symbolic coding.  The triangular Lorentz geometry is stated with the
correct strict corridor threshold.

\section*{Disclosure and review status}
The author is responsible for all analytical arguments, citations, and
computations.  Structural checks do not certify the proofs.  This revision is
prepared for a new independent formal review.

\bibliographystyle{plain}
\bibliography{references-round4}

\end{document}
